\chapter{Probability Models, Conditioning, and Random Variables}
\label{chap:probability-models}

Probability starts with a model: a precise description of what may happen,
which statements about the outcome are meaningful, and how probabilities are
assigned. Conditioning changes the reference population, while a random
variable turns outcomes into numerical values whose distribution can be
studied.

\section{Probability spaces and events}

\begin{definition}[Probability space]
A probability space is a triple
\[
  (\Omega,\Field,\Prob).
\]
Here \(\Omega\) is the \emph{sample space}, whose elements are outcomes;
\(\Field\) is a collection of subsets of \(\Omega\), called events; and
\(\Prob\) assigns a probability to every event.
\end{definition}

The event collection \(\Field\) is closed under complements and countable
unions:
\[
  A\in\Field \Longrightarrow A^c\in\Field,
  \qquad
  A_1,A_2,\ldots\in\Field
  \Longrightarrow \bigcup_{i=1}^{\infty}A_i\in\Field.
\]
It also contains \(\Omega\). In a finite model one normally takes
\(\Field=2^\Omega\), the set of all subsets of \(\Omega\).

For events \(A,B\subseteq\Omega\):
\begin{itemize}
  \item \(A\cap B\) means that both events occur;
  \item \(A\cup B\) means that at least one occurs;
  \item \(A^c=\Omega\setminus A\) means that \(A\) does not occur;
  \item \(A\setminus B=A\cap B^c\) means that \(A\) occurs but \(B\) does not.
\end{itemize}

\begin{definition}[Probability axioms]
For events in \(\Field\), a probability measure satisfies:
\begin{align*}
  &\Prob(A)\ge 0,\\
  &\Prob(\Omega)=1,\\
  &\Prob\!\left(\bigcup_{i=1}^{\infty}A_i\right)
    =\sum_{i=1}^{\infty}\Prob(A_i)
    \quad\text{for pairwise disjoint }A_i.
\end{align*}
\end{definition}

Several basic rules follow from these axioms:
\begin{align}
  \Prob(\varnothing)&=0,\label{eq:empty}\\
  \Prob(A^c)&=1-\Prob(A),\label{eq:complement}\\
  A\subseteq B&\Longrightarrow \Prob(A)\le \Prob(B),\label{eq:monotonicity}\\
  \Prob(A\cup B)&=\Prob(A)+\Prob(B)-\Prob(A\cap B).
  \label{eq:inclusion-exclusion}
\end{align}
The last identity is the inclusion--exclusion rule. Adding \(\Prob(A)\) and
\(\Prob(B)\) counts the overlap twice, so it must be subtracted once.

\begin{example}[Two fair coin tosses]
Toss a fair coin twice and record the ordered result. Then
\[
  \Omega=\{HH,HT,TH,TT\},
  \qquad \Prob(\{\omega\})=\frac14
  \quad\text{for every }\omega\in\Omega.
\]
Let
\[
  A=\{\text{first toss is }H\}=\{HH,HT\},
  \qquad
  B=\{\text{second toss is }H\}=\{HH,TH\}.
\]
Then
\[
  A\cap B=\{HH\},\qquad
  A\cup B=\{HH,HT,TH\},\qquad
  A^c=\{TH,TT\}.
\]
Inclusion--exclusion gives
\[
  \Prob(A\cup B)=\frac12+\frac12-\frac14=\frac34.
\]
\end{example}

For a finite sample space with equally likely outcomes,
\[
  \Prob(A)=\frac{|A|}{|\Omega|}.
\]
This is a counting shortcut, not an additional axiom: it is valid only after
equal likelihood has been justified.

\section{Conditional probability and total probability}

\begin{definition}[Conditional probability]
For \(\Prob(B)>0\), the probability of \(A\) given \(B\) is
\begin{equation}
  \Prob(A\mid B)=\frac{\Prob(A\cap B)}{\Prob(B)}.
  \label{eq:conditional-probability}
\end{equation}
\end{definition}

Conditioning restricts the reference set to \(B\) and renormalizes the
remaining probabilities. The event after the vertical bar determines the
denominator. In general, \(\Prob(A\mid B)\) and \(\Prob(B\mid A)\) are
different quantities.

Rearranging \eqref{eq:conditional-probability} gives the multiplication rule
\begin{equation}
  \Prob(A\cap B)=\Prob(A\mid B)\Prob(B)
  =\Prob(B\mid A)\Prob(A).
  \label{eq:multiplication-rule}
\end{equation}

\begin{example}[At least one head]
In the two-toss model, let
\[
  E=\{\text{at least one head}\}=\{HH,HT,TH\}.
\]
Given \(E\), the probability that the second toss is heads is
\[
  \Prob(B\mid E)
  =\frac{\Prob(B\cap E)}{\Prob(E)}
  =\frac{2/4}{3/4}=\frac23.
\]
The unconditional probability is \(\Prob(B)=1/2\). Conditioning removes the
outcome \(TT\), in which \(B\) fails.
\end{example}

\begin{theorem}[Law of total probability]
Let \(B_1,\ldots,B_m\) be a partition of \(\Omega\): the events are pairwise
disjoint and their union is \(\Omega\). If \(\Prob(B_i)>0\), then
\begin{equation}
  \Prob(A)=\sum_{i=1}^{m}\Prob(A\mid B_i)\Prob(B_i).
  \label{eq:total-probability}
\end{equation}
\end{theorem}

Indeed, the events \(A\cap B_i\) are disjoint and their union is \(A\). Add
their probabilities and then apply the multiplication rule. For a binary
partition \(B,B^c\),
\[
  \Prob(A)=\Prob(A\mid B)\Prob(B)
  +\Prob(A\mid B^c)\Prob(B^c).
\]

\section{Independence and mutual exclusivity}

\begin{definition}[Independence]
Events \(A\) and \(B\) are independent if
\begin{equation}
  \Prob(A\cap B)=\Prob(A)\Prob(B).
  \label{eq:independence}
\end{equation}
If \(\Prob(B)>0\), this is equivalent to
\(\Prob(A\mid B)=\Prob(A)\).
\end{definition}

Independence means that learning whether one event occurred does not change the
probability of the other. In the coin example,
\[
  \Prob(A\cap B)=\frac14=\frac12\cdot\frac12,
\]
so the events ``first toss is heads'' and ``second toss is heads'' are
independent.

\begin{definition}[Mutual exclusivity]
Events \(A\) and \(B\) are mutually exclusive, or disjoint, if
\[
  A\cap B=\varnothing.
\]
They cannot occur together.
\end{definition}

\begin{proposition}
Two mutually exclusive events with positive probability are not independent.
\end{proposition}

\begin{proof}
Mutual exclusivity gives \(\Prob(A\cap B)=0\), while positive probabilities
give \(\Prob(A)\Prob(B)>0\). Thus \eqref{eq:independence} cannot hold.
\end{proof}

The distinction is essential:
\begin{center}
\begin{tabular}{@{}L{0.28\textwidth}L{0.30\textwidth}L{0.30\textwidth}@{}}
\toprule
 & \textbf{Mutual exclusivity} & \textbf{Independence} \\
\midrule
Can both occur? & No & Usually yes \\
Defining condition & \(A\cap B=\varnothing\) &
\(\Prob(A\cap B)=\Prob(A)\Prob(B)\) \\
Information meaning & One rules out the other & One does not change the other \\
\bottomrule
\end{tabular}
\end{center}

\section{Bayes' theorem and base rates}

\begin{theorem}[Bayes' theorem]
Let \(B_1,\ldots,B_m\) be a partition of \(\Omega\). For \(\Prob(A)>0\),
\begin{equation}
  \Prob(B_j\mid A)
  =\frac{\Prob(A\mid B_j)\Prob(B_j)}
  {\sum_{i=1}^{m}\Prob(A\mid B_i)\Prob(B_i)}.
  \label{eq:bayes}
\end{equation}
\end{theorem}

The numerator is the probability of the route \(B_j\) followed by \(A\). The
denominator adds every mutually exclusive route that can produce \(A\).

\begin{example}[Rare-defect classifier]
Let \(D\) mean that an item is defective and \(+\) mean that a classifier
produces a positive flag. Assume
\[
  \Prob(D)=0.01,\qquad
  \Prob(+\mid D)=0.99,\qquad
  \Prob(+\mid D^c)=0.05.
\]
The probability of a positive result is
\[
  \Prob(+)=0.99(0.01)+0.05(0.99)=0.0594.
\]
Therefore
\[
  \Prob(D\mid +)
  =\frac{0.99(0.01)}{0.0594}
  =\frac16\approx 0.167.
\]

The same calculation can be expressed with expected counts in a population of
\(10{,}000\) items:
\begin{center}
\begin{tabular}{@{}lrrr@{}}
\toprule
\textbf{Actual state} & \textbf{Positive} & \textbf{Negative} & \textbf{Total} \\
\midrule
Defect \(D\) & 99 & 1 & 100 \\
No defect \(D^c\) & 495 & 9,405 & 9,900 \\
\midrule
Total & 594 & 9,406 & 10,000 \\
\bottomrule
\end{tabular}
\end{center}
Only \(99\) of the \(594\) positive items are defective. The sensitivity
\(\Prob(+\mid D)=0.99\) is not the posterior \(\Prob(D\mid +)\). The low base
rate and the false positives among the much larger non-defective group dominate
the denominator.
\end{example}

For a Bayesian classifier with classes \(C_1,\ldots,C_K\) and observed features
\(x\),
\begin{equation}
  \Prob(C_k\mid x)
  =\frac{p(x\mid C_k)\Prob(C_k)}
  {\sum_{j=1}^{K}p(x\mid C_j)\Prob(C_j)}.
  \label{eq:bayesian-classifier}
\end{equation}
The factors have distinct roles:
\begin{itemize}
  \item \(\Prob(C_k)\) is the prior or class base rate;
  \item \(p(x\mid C_k)\) is the class-conditional likelihood;
  \item \(\Prob(C_k\mid x)\) is the posterior after observing \(x\).
\end{itemize}
The maximum a posteriori rule is
\[
  \widehat C_{\mathrm{MAP}}(x)
  =\argmax_k p(x\mid C_k)\Prob(C_k).
\]

\section{Random variables}

\begin{definition}[Random variable]
A real-valued random variable is a function
\[
  X:\Omega\longrightarrow\Real
\]
such that \(\{\omega:X(\omega)\le x\}\) is an event for every
\(x\in\Real\). In finite examples this event condition is automatic.
\end{definition}

The sampled outcome is random; the mapping \(X\) is fixed.

\begin{example}[Number of heads]
For two coin tosses, let \(X\) be the number of heads. Then
\begin{center}
\begin{tabular}{@{}ccccc@{}}
\toprule
Outcome \(\omega\) & \(HH\) & \(HT\) & \(TH\) & \(TT\) \\
\midrule
\(X(\omega)\) & 2 & 1 & 1 & 0 \\
\bottomrule
\end{tabular}
\end{center}
The support \(\{0,1,2\}\) is the set of possible values of \(X\); it is not the
original sample space.
\end{example}

A \emph{discrete} random variable takes values in a finite or countable set. An
\emph{absolutely continuous} random variable assigns probabilities to intervals
through a density. Mixed random variables, with both point masses and a
continuous part, also exist.

\section{Exercises}

\begin{exercise}
For a fair die, let \(A=\{2,4,6\}\) and \(B=\{4,5,6\}\). Find
\(A\cap B\), \(A\cup B\), \(A^c\), their probabilities, and
\(\Prob(A\mid B)\). Are \(A\) and \(B\) independent?
\end{exercise}

\begin{exercise}
A system uses server type \(A\) with probability \(0.6\) and type \(B\) with
probability \(0.4\). Their failure probabilities are \(0.02\) and \(0.05\),
respectively. Find the total failure probability and the posterior probability
that a failed server was type \(B\).
\end{exercise}

\begin{exercise}
For a fair die, show that \(A=\{2,4,6\}\) and \(B=\{1,2,3,4\}\) are
independent. Compare this with \(A\) and \(C=\{1,3,5\}\).
\end{exercise}

\begin{exercise}
A binary classifier has classes \(C_0,C_1\), with \(\Prob(C_1)=0.2\),
\(\Prob(s\mid C_1)=0.6\), and \(\Prob(s\mid C_0)=0.1\). Compute
\(\Prob(C_1\mid s)\) and select the MAP class.
\end{exercise}

\subsection*{Short answers}

\begin{enumerate}
  \item \(A\cap B=\{4,6\}\), \(A\cup B=\{2,4,5,6\}\), and
  \(A^c=\{1,3,5\}\). Their probabilities are \(1/3\), \(2/3\), and
  \(1/2\), respectively; \(\Prob(A\mid B)=2/3\). Since
  \(\Prob(A\cap B)=1/3\ne\Prob(A)\Prob(B)=1/4\), the events are not
  independent.
  \item \(\Prob(F)=0.032\) and \(\Prob(B\mid F)=0.625\).
  \item \(\Prob(A\cap B)=1/3=\Prob(A)\Prob(B)\); \(A\) and \(C\) are
  mutually exclusive and not independent.
  \item \(\Prob(C_1\mid s)=0.6\), so the MAP class is \(C_1\).
\end{enumerate}
